\subsubsection{Compactness Properties}\label{subsec:sobolev-compactness-properties}

In this subsection, we show that certain imbeddings involving Sobolev spaces are compact. This is often useful in applications since it implies that certain sequences have convergent subsequences; the limit of such a subsequence is typically the solution that is being sought. It is also useful in defining certain equivalent norms for Sobolev spaces as we do at the end of this subsection.

We begin with the definition of a compact imbedding.

\begin{Definition}
\label{def:compact-imbedding}
Let $X,Y$ be Banach spaces such that $X$ is continuously imbedded in $Y$. We say that $X$ is \textbf{compactly imbedded} in $Y$ and write\[X\overset{c}{\hookrightarrow}Y\]if the unit ball in $X$ is precompact in $Y$ or, equivalently, every bounded sequence in $X$ has a subsequence that converges in $Y$.
\end{Definition}

Before identifying compact imbeddings of Sobolev spaces, we establish a lemma.

\begin{Lemma}
\label{lem:sobolev-holder-continuity}
Let $s>m/2$ and $0<\alpha<\min(s-m/2,1)$. Then there is a constant $C$ such that
\begin{equation}\label{eq:sobolev-holder-continuity-estimate}
|u(\bm{x})-u(\bm{y})|\leq C\|u\|_{s,2}|\bm{x}-\bm{y}|^{\alpha}
\end{equation}
for every $\bm{x},\bm{y}\in\mathbb{R}^m$ and $u\in H^s(\mathbb{R}^m)$.
\end{Lemma}

\begingroup
\postdisplaypenalty=10000
\renewcommand{\endProof}{\par\nobreak\hfill$\blacksquare$\par\vspace{0.1cm}}
\begin{Proof}
Throughout the proof, let $C$ denote a generic constant independent of $\bm{x}$, $\bm{y}$ and $u$. We first observe that, for any $\alpha\in(0,1)$,
\begin{equation}\label{eq:fourier-exponential-holder-difference}
|\exp(i\bm{\xi}\bm{x})-\exp(i\bm{\xi}\bm{y})|\leq C|\bm{x}-\bm{y}|^{\alpha}|\bm{\xi}|^{\alpha}.
\end{equation}
With $\hat u$ denoting the Fourier transform of $u$, we have
\begin{equation}\label{eq:sobolev-holder-fourier-estimate}
\begin{aligned}
|u(\bm{x})-u(\bm{y})|&\leq(2\pi)^{-m/2}\int_{\mathbb{R}^m}|\exp(i\bm{\xi}\bm{x})-\exp(i\bm{\xi}\bm{y})||\hat u(\bm{\xi})|\,\mathrm{d}\bm{\xi}\\
&\leq C|\bm{x}-\bm{y}|^{\alpha}\int_{\mathbb{R}^m}|\hat u(\bm{\xi})||\bm{\xi}|^{\alpha}\,\mathrm{d}\bm{\xi}\\
&\leq C|\bm{x}-\bm{y}|^{\alpha}\|(1+|\bm{\xi}|^s)\hat u\|_2\||\bm{\xi}|^{\alpha}/(1+|\bm{\xi}|^s)\|_2\\
&\leq C|\bm{x}-\bm{y}|^{\alpha}\|u\|_{s,2}.
\end{aligned}
\end{equation}

\end{Proof}
\endgroup

We now prove one of the most basic compact imbedding theorems.

\begin{Theorem}
\label{thm:sobolev-continuous-target-compact-embedding}
Let $k>m/2$ and let $\Omega$ be bounded and such that it has the $k$-extension property. Then\[H^k(\Omega)\overset{c}{\hookrightarrow}C_b(\overline\Omega),\]i.e., $H^k(\Omega)$ is compactly imbedded in $C_b(\overline\Omega)$.
\end{Theorem}

\begin{Proof}
The theorem follows immediately from the Sobolev imbedding theorem (Theorem~\ref{thm:sobolev-full-space-continuous-embedding}), Lemma~\ref{lem:sobolev-holder-continuity}, and the Arzela-Ascoli theorem (Theorem~\ref{thm:arzela-ascoli}).
\end{Proof}

\Needspace{4\baselineskip}
\begin{Example}
\label{ex:unbounded-domain-translated-functions-noncompactness}
The assumption that $\Omega$ is bounded is essential. To see this let $\Omega=\mathbb{R}^m$ and consider the sequence of translated functions\[u_n(\bm{x})=u(\bm{x}+n\bm{e}),\]where $\bm{e}$ is any fixed unit vector and $u$ any nonzero element of $H^k(\mathbb{R}^m)$. This sequence is, of course, bounded in $H^k(\mathbb{R}^m)$. However, the sequence can have no uniformly convergent subsequence. (To prove this, one needs to show that no such sequence that does converge uniformly can be in $H^k(\mathbb{R}^m)$ because of its behavior at infinity.)
\end{Example}
